\chapter{Em Có Biết Gì Đâu}
\chapterintro{Chủ đề nối thêm}{Chương này chuyển hẳn sang giọng thầy cô ghẹo học trò bằng những câu ngắn. Mọi bài đều được căn chỉnh tuyệt đối theo luật bằng trắc ngũ ngôn (vần bằng, nhịp Trắc-Bằng luân phiên) để vừa vui vừa thật chuẩn nhịp.}

\poem{Bài 101}{Ngũ ngôn tứ tuyệt}{Em có biết gì đâu}{Bài thầy giao đã lâu}{Nghe kêu mồm ấp úng}{Vuốt trán xoa vò đầu}
\poem{Bài 102}{Ngũ ngôn tứ tuyệt}{Em có biết gì đâu}{Mơ màng say giấc sâu}{Cô hô liền bật dậy}{Đánh rớt bay từng câu}
\poem{Bài 103}{Ngũ ngôn tứ tuyệt}{Em có biết gì đâu}{Cùng ngồi sao chép nhau}{Sai ly sai dấu phẩy}{Chỉ biết ôm niềm sầu}
\poem{Bài 104}{Ngũ ngôn tứ tuyệt}{Em có biết gì đâu}{Tuần này quên tiết đầu}{Ba lô mang tập trắng}{Chống chế bao lời ngầu}
\poem{Bài 105}{Ngũ ngôn tứ tuyệt}{Em có biết gì đâu}{Buôn lời trôi hết sầu}{Thầy kêu lên hỏi khó}{Bí quá đôi tròng ngầu}
\poem{Bài 106}{Ngũ ngôn tứ tuyệt}{Em có biết gì đâu}{Bút này phai đã lâu}{Chưa ghi đành nín lặng}{Báo cáo rưng rưng sầu}
\poem{Bài 107}{Ngũ ngôn tứ tuyệt}{Em có biết gì đâu}{Bài văn đi quá sâu}{Cô phê dòng chói lọi}{Cúi mặt ôm nghìn sầu}
\poem{Bài 108}{Ngũ ngôn tứ tuyệt}{Em có biết gì đâu}{Đường dài kẹt giữa cầu}{Trà chanh tay vẫn ấm}{Cả lớp cười xoa đầu}
\poem{Bài 109}{Ngũ ngôn tứ tuyệt}{Em có biết gì đâu}{Ra ngoài sao quá lâu}{Hồi chuông vang hết tiết}{Mới dám thò nhô đầu}
\poem{Bài 110}{Ngũ ngôn tứ tuyệt}{Em có biết gì đâu}{Lỗi này quen đã lâu}{Ngày mai thề sửa đổi}{Tối lại y như đầu}
